\section{Knowledge of the Heart}

Light is the most common symbol of knowledge, so that it is natural that solar light represent direct knowledge, that is, intuitive knowledge, which is that of pure intellect, and that lunar light represent reflective knowledge, that is, discursive knowledge, which is that of reason. As the moon can only give its light when it is itself illuminated by the sun, so likewise reason, within the order of reality that is its own proper domain, can function validly only under the guarantee of principles which enlighten and direct it, and which it receives from the higher intellect. On this point there is an ambiguity which it is important to dispel: modern philosophers strangely deceive themselves in speaking as they do of ``rational principles", as if these principles belonged by right to reason, as if they were in some way its production, whereas in order to govern reason they must on the contrary impose themselves upon it, and thus come from above. This is an example of the rationalist error, and one can thereby understand the essential difference that exists between rationalism and true intellectuality. One need reflect but a moment to understand that a principle in the true sense of the word, by the very fact that it cannot be drawn or deduced from something else, can only be grasped immediately, thus intuitively, and could not be the object of a discursive knowledge such as that which characterizes reason to use Scholastic terminology, it is pure intellect that is \emph{habitus principorum}, whereas reason is only \emph{habitus conclusionum}.

A further consequence results from the fundamental character of intellect and reason: intuitive knowledge, because it is immediate, is necessarily infallible in itself; on the contrary, error can always creep into any knowledge that is only indirect or mediate, as is rational knowledge, and it can be seen thereby how wrong Descartes was in wishing to attribute infallibility to reason. This is what \textbf{Aristotle} expresses in these terms:

\begin{quotex}
Among the properties of intelligence, in virtue of which we attain to truth, there are some that are always true, and others that may lead to error. Reason is in the latter class; but the intellect is always in conformity with truth, and nothing is truer than the intellect. Now, principles being better known than demonstration, and every science being accompanied by reason, the knowledge of principles is not a science [but it is a mode of knowledge superior to scientific or rational knowledge, and constitutes metaphysical knowledge proper.] Besides, only the intellect is truer than science [or than reason, which builds up science]; thus principles derive from the intellect. 

\end{quotex}
And to better affirm the intuitive character of this intellect, Aristotle goes on to say:

\begin{quotex}
One does not demonstrate principles, but one perceives their truth directly. 

\end{quotex}
This direct perception of truth, this intellectual and supra-rational intuition, the simplest notion of which modern man seems to have lost, is true ``knowledge of the heart", to use an expression frequently found in Eastern doctrines. Moreover, this knowledge is in itself something incommunicable; it is necessary to have ``realized" it, at least in a certain measure, to know what it truly is, and all that can be said about it gives only la more or less approximate but always inadequate idea of it. Above all, it would be an error to believe that the nature of such knowledge could be effectively understood by anyone who is content to look at it ``philosophically", that is, from outside, for it must never be forgotten that like all ``profane learning" philosophy is only a purely human or rational knowledge.

On the contrary, ``sacred science", in the sense that we have used this term in our writings, is essentially based on supra-rational knowledge; and all that we have said about the use of symbolism and of the teaching contained therein is related to the means which the traditional doctrines place at the disposal of man to enable him to attain to that knowledge par excellence with regard to which all other knowledge, in the measure that it too has some reality, is only a more or less remote participation, a more or less indirect reflection, just as the light of the moon is only a pale reflection of that of the sun.

``Knowledge of the Heart" is the direct perception of the intelligible light, of that Light of the Word spoken of by St. John in the prologue of his Gospel, that radiant Light of the ``Spiritual Sun" which is the true ``Heart of the World".

\flright{\textsc{René Guénon}, Heart and Brain, from \textit{Symbols of Sacred Science}}

\flrightit{Posted on 2006-10-25 by Cologero }
